\section*{Chapter \ref{ch:b1:consumer-welfare} --- Consumer Welfare}
% Numbers recomputed in code/bachelor-1/07-consumer-welfare/exercises/solutions.py (tested).

\begin{solution}{exo:b1:consumer-welfare:cs-linear}
Quantities 80 at a price of 10 and 70 at 15. The surplus lost is the
trapezium between the two prices: $5 \times (80 + 70)/2 = 375$.
\end{solution}

\begin{solution}{exo:b1:consumer-welfare:growth}
CPI-U: $2.4001^{1/34} - 1 = 2.61\%$ a year. Deflator: $2.0689^{1/34} - 1 = 2.16\%$
a year.
\end{solution}

\begin{solution}{exo:b1:consumer-welfare:lp}
Laspeyres: $(2 \times 10 + 1 \times 10)/(1 \times 10 + 1 \times 10) = 30/20 = 1.5$.
Paasche: $(2 \times 6 + 1 \times 12)/(1 \times 6 + 1 \times 12) = 24/18 = 1.333$.
\end{solution}

\begin{solution}{exo:b1:consumer-welfare:signs}
Both are negative. $\mathrm{CV} < 0$: at the new, lower price the consumer
could give up $|\mathrm{CV}|$ of income and be as well off as before.
$\mathrm{EV} < 0$: at the old price she would need $|\mathrm{EV}|$ more income to
be as well off as the price fall makes her.
\end{solution}

\begin{solution}{exo:b1:consumer-welfare:cv-ev}
$e(p, u) = 2u\sqrt{p_1p_2}$, so doubling $p_1$ multiplies the cost of any
utility by $\sqrt 2$. $\mathrm{CV} = e(p^1, u^0) - 1000 = 1000(\sqrt 2 - 1) = 414.21$.
$\mathrm{EV} = 1000 - e(p^0, u^1) = 1000 - 1000/\sqrt 2 = 292.89$.
\end{solution}

\begin{solution}{exo:b1:consumer-welfare:between}
$x_1 = 500/p_1$, so the loss is $\int_1^2 500/p\,\dd p = 500\ln 2 = 346.57$, and
$292.89 < 346.57 < 414.21$.
\end{solution}

\begin{solution}{exo:b1:consumer-welfare:cd-index}
With $u = \prod_i x_i^{s_i}$, the \omterm{def:b1:income-and-substitution-effects:expenditure}{expenditure function} is
$e(p, u) = u\prod_i p_i^{s_i}/\prod_i s_i^{s_i}$ (solve $v(p, m) = u$ for $m$, as in
\cref{ch:b1:income-and-substitution-effects}). In the ratio $e(p^1, u)/e(p^0, u)$
the factor $u/\prod_i s_i^{s_i}$ cancels, leaving $\prod_i(p_i^1/p_i^0)^{s_i}$ for
every $u$. For the figure's consumer, shares $\tfrac12$ and relatives 2 and 1:
$I = 2^{1/2} = 1.4142$, against $L = (2 \times 50 + 1 \times 50)/100 = 1.5$: an error of
$0.0858$, or 8.6 index points.
\end{solution}

\begin{solution}{exo:b1:consumer-welfare:basket}
Laspeyres: $(600 \times 1.355 + 800 \times 1.225 + 750 \times 1.84 + 8400 \times 1.4288 + 300 \times 1.3033)/10850
= 15565.68/10850 = 1.4346$, that is 143.46 (the products are the 2025 costs of the
basket, 813, 980, 1380, 12001.68 and 391, using the 2025 prices). Geometric:
the share-weighted logs are $0.0168 + 0.0150 + 0.0421 + 0.2763 + 0.0073 = 0.3575$,
and $100e^{0.3575} = 142.98$ with these rounded terms (142.97 unrounded, as in
the table). The chained index, 143.27, lies between them, closer to the
Laspeyres index.
\end{solution}

\begin{solution}{exo:b1:consumer-welfare:spread}
$r_i - \bar r = \pm 0.01$ in one year, so the bias is
$\tfrac12(0.5 \times 0.0001 + 0.5 \times 0.0001) = 0.00005$: $0.005$ points. A basket fixed
for ten years faces relatives of $\pm 0.10$, so the bias in the level is
$\tfrac12 \times 0.01 = 0.005$, half a per cent after ten years, about $0.05$
points a year: a fixed basket's bias grows with the square of the time
since its base year.
\end{solution}

\begin{solution}{exo:b1:consumer-welfare:prove-bound}
The proof is that of \cref{prop:b1:consumer-welfare:bounds}: $x^0$ reaches $u^0$,
so $e(p^1, u^0) \le p^1\cdot x^0$ while $e(p^0, u^0) = p^0\cdot x^0$; and
$e(p^0, u^1) \le p^0\cdot x^1$ while $e(p^1, u^1) = p^1\cdot x^1$. With
$u = \min(x_1/a, x_2/b)$, the cheapest way to reach $u$ is $(au, bu)$ at any
prices, so $e(p, u) = u(ap_1 + bp_2)$ and the basket never changes:
$I = (ap_1^1 + bp_2^1)/(ap_1^0 + bp_2^0) = L$. With no substitution there is no
substitution bias.
\end{solution}

\begin{solution}{exo:b1:consumer-welfare:quasi}
$5/\sqrt{x_1} = p_1$ gives $x_1 = 25/p_1^2$: 25 at $p_1 = 1$ (cost 25) and 6.25 at
$p_1 = 2$ (cost 12.5), both interior. The surplus loss is
$\int_1^2 25/p^2\,\dd p = 25(1 - \tfrac12) = 12.5$. The \omterm{def:b1:income-and-substitution-effects:expenditure}{expenditure function} is
$e(p, u) = p_1x_1 + u - 10\sqrt{x_1} = 25/p_1 + u - 50/p_1 = u - 25/p_1$, so
$e(2, u) - e(1, u) = -12.5 + 25 = 12.5$ for every $u$: $\mathrm{CV} = \mathrm{EV} = 12.5$.
\end{solution}

\begin{solution}{exo:b1:consumer-welfare:low-income}
A price index built on the baskets of low-income households, with their own
shares, that rose faster than the CPI by more than the CPI's estimated
upward bias would refute the statement for them. Their true index can rise
faster if they spend larger shares on goods whose prices rose faster than
average (rent, food, energy in some periods), or if they can substitute
less, for instance because they cannot buy in bulk or travel to cheaper
outlets.
\end{solution}

\begin{solution}{pb:b1:consumer-welfare:cpi}
\textbf{1.} $600 + 800 + 750 + 8400 + 300 = 10850$. Rent $8400/10850 = 0.7742$;
energy $750/10850 = 0.0691$.

\textbf{2.} $300 \times 2.71 + 20 \times 49 + 3000 \times 0.46 + 12 \times 1000.14 + 100 \times 3.91 = 813 + 980 + 1380 + 12001.68 + 391 = 15565.68$;
$L = 100 \times 15565.68/10850 = 143.46$.

\textbf{3.} $p^1\cdot x^0/p^0\cdot x^0 = \sum_i (p_i^0x_i^0/p^0\cdot x^0)(p_i^1/p_i^0) = \sum_i s_ir_i$
with $r_i = p_i^1/p_i^0$.

\textbf{4.} Cobb--Douglas demand spends the share $s_i$ of income $m^1$ on good $i$:
$x_i^1 = s_im^1/p_i^1$. Then $p^1\cdot x^1 = m^1$ and
$p^0\cdot x^1 = m^1\sum_i s_ip_i^0/p_i^1$, so $P = 1/\sum_i s_i/r_i$.

\textbf{5.} $\sum_i s_i/r_i = (600/1.355 + 800/1.225 + 750/1.84 + 8400/1.4288 + 300/1.3033)/10850
= (442.80 + 653.06 + 407.61 + 5879.06 + 230.18)/10850 = 7612.71/10850$, so
$P = 100 \times 10850/7612.71 = 142.52$.

\textbf{6.} $\ln G = \sum_i s_i\ln r_i = 0.0168 + 0.0150 + 0.0421 + 0.2763 + 0.0073 = 0.3575$
and $G = 100e^{0.3575} = 142.98$ with these rounded terms; the script's
unrounded value is 142.97.

\textbf{7.} $L = 143.46 > G = 142.97 > P = 142.52$: the share-weighted arithmetic,
geometric and harmonic means of the price relatives, in the order that holds
for any set of unequal numbers.

\textbf{8.} The Cobb--Douglas \omterm{def:b1:income-and-substitution-effects:expenditure}{expenditure function} is proportional to
$\prod_i p_i^{s_i}$ at every utility, so $e(p^1, u)/e(p^0, u) = \prod_i r_i^{s_i} = G$
for every reference utility (\cref{exo:b1:consumer-welfare:cd-index}).

\textbf{9.} $143.46 - 142.97 = 0.49$ index points.

\textbf{10.} $1.4346^{0.1} - 1 = 3.67\%$ and $1.4297^{0.1} - 1 = 3.64\%$: $0.03$ points
a year with these rounded rates ($0.035$ unrounded), under $0.04$ points a year.

\textbf{11.} Rent takes $77\%$ of the basket and its relative, 1.4288, is close to
the mean; substitution matters only when relative prices move apart, and
only energy, with $7\%$ of the basket, moved far (1.84).

\textbf{12.} Energy prices jumped in 2022. The relatives spread out, and the
bias of a fixed basket grows with the spread
(\cref{exo:b1:consumer-welfare:spread}).

\textbf{13.} The chained index, 143.27, lies between the Laspeyres and the
geometric index. Each yearly link is still an arithmetic mean of that year's
relatives, so it carries a small bias every year, and the 2022 link a large
one; chaining stops the bias from compounding on an old basket but does not
remove it.

\textbf{14.} 30 dalens per gigabyte, where $60 - 2p = 0$.

\textbf{15.} 2020: $q = 20$, $\mathrm{CS} = (30 - 20) \times 20/2 = 100$. 2025: $q = 40$,
$\mathrm{CS} = (30 - 10) \times 40/2 = 400$ dalens a year.

\textbf{16.} With \omterm{def:b1:preferences-and-utility:quasi-linear}{quasi-linear utility} the surplus is exact
(\cref{prop:b1:consumer-welfare:quasi}): reaching the 2015 utility costs 400
dalens less, $400/15565.68 = 2.57\%$ of the 2025 cost of the 2015 basket.

\textbf{17.} $10850 \times 142.97/100 = 15512.25$; with the phone,
$15512.25 - 400 = 15112.25$, an index of $100 \times 15112.25/10850 = 139.28$.

\textbf{18.} $1.3928^{0.1} - 1 = 3.37\%$ against $3.64\%$: an index that leaves the
phone out overstates by $0.27$ points a year, several times the substitution
bias.

\textbf{19.} The whole value of the new good: the price fall from the \omterm{def:b1:consumer-welfare:virtual-price}{virtual
price}, 30, to 20 when it arrived, and from 20 to 10 since. Adding it in 2025
at 10 records neither.

\textbf{20.} $15000 \times 1.4346 = 21519$ and $15000 \times 1.4297 = 21445.50$.

\textbf{21.} $21519 - 21445.50 = 73.50$ dalens in 2025, and
$15000 \times 1.81/100 = 271.50$ dalens over 2016--2025.

\textbf{22.} $15000 \times 1.4327 = 21490.50$; the switch would cost her
$21519 - 21490.50 = 28.50$ dalens in 2025.

\textbf{23.} The geometric index, 142.97: it is the true index for a household
with these preferences and shares that does not use the new good.

\textbf{24.} It understates hers. Energy's relative, 1.84, is the largest, and her
true index weights it more, so her geometric index is higher than the
average one.

\textbf{25.} The substitution bias of Brenland's CPI over 2015--2025 is 0.49
index points in all, under $0.04$ points a year; leaving out the new good
costs far more, $0.27$ points a year.
\end{solution}
